\documentclass[11pt,a4paper]{article}
\usepackage[T1]{fontenc}
\usepackage[utf8]{inputenc}
\usepackage{lmodern}
\usepackage{amsmath,amssymb,amsthm,mathtools,mathrsfs}
\usepackage{graphicx,float,xcolor,multicol}
\usepackage[shortlabels]{enumitem}
\usepackage[margin=27mm]{geometry}
\usepackage{microtype,needspace,placeins}
\usepackage{tikz,tikz-cd}
\usetikzlibrary{arrows.meta,calc,positioning,patterns,decorations.pathreplacing}
\usepackage[colorlinks=true,linkcolor=teal,urlcolor=teal,citecolor=teal]{hyperref}
\setlength{\parindent}{0pt}
\setlength{\parskip}{.6em}
\setcounter{secnumdepth}{0}
\allowdisplaybreaks
\raggedbottom
\providecommand{\cross}{\times}
\DeclareUnicodeCharacter{2020}{\ensuremath{\dagger}}
\DeclareMathOperator{\supp}{supp}
\DeclareMathOperator*{\esssup}{ess\,sup}
\providecommand{\chapter}{\section}
\newenvironment{tool}[2]{\subsection*{Tool #1 --- #2}}{}
\newcommand{\ToolRef}[1]{Tool #1}
\newcommand{\FigureTag}[2]{\textit{Figure #1}}
\newcommand{\Result}[1]{\subsection*{#1}}
\newcommand{\Application}[1]{\subsubsection*{Application\if\relax\detokenize{#1}\relax\else\space--- #1\fi}}
\newcommand{\ProofHeading}{\par\textit{Proof.}\space}
\newcommand{\ProofOf}[1]{\par\textit{Proof of #1.}\space}
\newcommand{\RemarkHeading}{\par\textbf{Remark.}\space}
\newcommand{\NoteHeading}{\par\textbf{Note.}\space}
\newcommand{\ExerciseHeading}{\par\textbf{Exercise.}\space}

\graphicspath{{figures/complex-analysis/}}
\begin{document}
\section*{Continuous Induction}
% BLOG-CONTENT-BEGIN
\subsection*{Tool 7 — Continuous Induction}
Recall strong induction. If a property $P$ is true for $n=1$, and if one can prove that
\[
P\text{ holds for every }k\leq n
\Longrightarrow P\text{ holds for }n+1,
\]
then $P$ holds for $\mathbb N$. There is a way to extend this to an interval $[0,1]$. Naturally, we expect $P$ to hold at $0$. There are several difficulties in devising an ``induction'' hypothesis for an interval.


\paragraph{Problem 1.}
If you demand the induction hypothesis to be
\[
P\text{ holds for every }t'\leq t
\Longrightarrow P\text{ holds for }t,
\]
then the induction stops at $0$.

\paragraph{Problem 2.}
If you demand that, whenever $P$ holds for $t\in[0,1]$, it holds for all \(t'\in(t,t+\delta)\)
for some $\delta$ depending on $t$, then the induction starts. But once $P$ holds on $[0,\delta)$, what if some $t\in[0,\delta)$ has $\delta_t$ such that
\[
(t,t+\delta_t)\subset[0,\delta)?
\]
The induction would stop there.


\paragraph{Problem 3.}
Let us demand both:
\[
P\text{ holds for every }t'<t
\Longrightarrow P\text{ holds for }t,
\]
and
\[
P\text{ holds for }t
\Longrightarrow
P\text{ holds for every }t'\in(t,t+\delta_t).
\]
Then too we have a problem. The induction progresses in this case, but it may not cover $[0,1]$. Start with $0$. The property holds for $0$, hence on $[0,\delta_0)$. It holds for every $t<\delta_0$, hence it holds for $\delta_0$. It then holds in a $\delta_0$-neighborhood of $\delta_0$, and then near the next endpoint, and so on. But the resulting sequence may converge to a point less than $1$. In fact, it may converge to a point arbitrarily close to $0$, in which case we have hardly covered $[0,1]$.

\paragraph{Problem 4.}
Let us say that the limit point is $\delta_1$. Since the assertion is true for every $\delta<\delta_1$, it holds for $\delta_1$. Start the induction again. We may end up with another limit point $\delta_2$. What if $\delta_2$ is also very close to $0$? In fact, what if the limit point of the sequence $\{\delta_n\}_n$ is very close to $0$?

\paragraph{Problem 5.}
What if we require the induction process never to stop? At a limit point, we would continue from that point, producing iterated limit points. The remaining task is to express this continuation precisely.


What do we do? We observe that $[0,1]$ is connected, so first demand that the set of points for which $P$ holds be closed in $[0,1]$. Then proceed as suggested: $P$ holds for $0$, and if $P$ holds for $t\in[0,1]$, then there is a $\delta_t$ such that every point of
\[
(t-\delta_t,t+\delta_t)
\]
also satisfies $P$. Thus the set for which $P$ holds is both open and closed, so $P$ holds on $[0,1]$. This is called the \emph{continuity method}, or continuous induction.

Let us state it differently. Let \(A:=\{t: \text{for every }t'\leq t,\ P\text{ holds for }t'\}.\)
Assume:
\begin{itemize}
\item $P$ holds for $0$;
\item $A$ is closed;
\item if $P$ holds for $t\in A$, then $P$ holds on $(t,t+\delta_t)$.
\end{itemize}
Closedness ensures that every limit point remains in the set, so the process can continue through limit points. Then $P$ holds on $[0,1]$.

This method is used to prove that $[0,1]$ is compact, which leads to the Heine--Borel theorem. Let us see two other instances of this method.

\subsection*{Arc-length formula}
If \(\gamma\colon[a,b]\to\mathbb C\)
is $C^1$ and rectifiable, then
\[
|\gamma|=\int_a^b|\gamma'(t)|\,dt.
\]

\textit{Proof.}
First note that
\[
|\gamma(t_j)-\gamma(t_{j-1})|
=\left|\int_{t_{j-1}}^{t_j}\gamma'(t)\,dt\right|.
\]
Let \(a=t_0<t_1<\cdots<t_n=b\)
be a partition. Then


\[
\sum_{j=1}^{n}|\gamma(t_j)-\gamma(t_{j-1})|
\leq\int_a^b|\gamma'(t)|\,dt,
\]
because
\[
\left|\int_a^b\gamma'(t)\,dt\right|
\leq\int_a^b|\gamma'(t)|\,dt.
\]
Taking the supremum over partitions gives
\[
|\gamma|\leq\int_a^b|\gamma'(t)|\,dt.
\]

It remains to prove
\[
|\gamma|\geq\int_a^b|\gamma'(t)|\,dt.
\]
For the other direction, it is enough to argue that, given any $\varepsilon>0$, there is a partition $P$ such that
\[
S_P(\gamma)>|\gamma|-\varepsilon.
\]
Here we use the same strategy: weaken the target inequality and prove
\[
|\gamma|\geq\int_a^b\bigl(|\gamma'(t)|-\varepsilon\bigr)\,dt.
\]
Let us prove this weak version inductively. Denote
\[
\left|\gamma\big|_{[a,t]}\right|=|\gamma|_t.
\]
At $a$,
\[
|\gamma|_a\geq\int_a^a(|\gamma'(t)|-\varepsilon)\,dt.
\]
Let
\[
\Omega_\varepsilon=
\left\{t:
\text{for every }t'\leq t,
\ |\gamma|_{t'}\geq
\int_a^{t'}(|\gamma'(s)|-\varepsilon)\,ds
\right\}.
\]
To check that $\Omega_\varepsilon$ is closed, it is enough to check whether
\[
\sup_{t\in\Omega_\varepsilon}t\in\Omega_\varepsilon.
\]
Let
\[
t_0=\sup_{t\in\Omega_\varepsilon}t.
\]
Since
\[
|\gamma|_{t_0}\geq|\gamma|_t
\qquad(t<t_0,\ t\in\Omega_\varepsilon),
\]
we get
\[
|\gamma|_{t_0}
\geq\sup_t\left\{
\int_a^t(|\gamma'(s)|-\varepsilon)\,ds
\right\}
=\int_a^{t_0}(|\gamma'(s)|-\varepsilon)\,ds.
\]


Thus $t_0\in\Omega_\varepsilon$, and $\Omega_\varepsilon$ is closed. We have one last thing to prove: there is a $\delta_t>0$ such that, for every
\[
t'\in[t,t+\delta_t),
\]
the inequality holds, whenever $t\in\Omega_\varepsilon$.

Since we know
\[
|\gamma|_t\geq\int_a^t(|\gamma'(s)|-\varepsilon)\,ds,
\]
we only have to prove
\[
\left|\gamma\big|_{[t,t']}\right|
\geq\int_t^{t'}(|\gamma'(s)|-\varepsilon)\,ds.
\]
Use the uniform continuity of $|\gamma'(t)|$. For $|t-t''|<\delta$, we have \(|\gamma'(t'')|\leq|\gamma'(t)|+\varepsilon/2.\)
Therefore, if $\delta_t<\delta$,
\[
\int_t^{t'}|\gamma'(s)|\,ds
\leq|\gamma'(t)|(t'-t)+\frac\varepsilon2(t'-t).
\]
If $\delta_t$ is small enough, then
\[
|\gamma'(t)|(t'-t)
\leq|\gamma(t')-\gamma(t)|+\frac\varepsilon2(t'-t),
\]
because
\[
\left|\frac{\gamma(t')-\gamma(t)}{t'-t}-\gamma'(t)\right|
<\frac\varepsilon2
\]
when $\delta_t$ is small enough. Consequently,
\[
|\gamma(t')-\gamma(t)|
\geq\int_t^{t'}(|\gamma'(s)|-\varepsilon)\,ds,
\]
and hence
\[
\left|\gamma\big|_{[t,t']}\right|
\geq\int_t^{t'}(|\gamma'(s)|-\varepsilon)\,ds.
\]
Continuous induction gives
\[
|\gamma|\geq\int_a^b(|\gamma'(t)|-\varepsilon)\,dt
\qquad(\forall\varepsilon>0),
\]
so
\[
|\gamma|\geq\int_a^b|\gamma'(t)|\,dt.
\]
Therefore
\[
|\gamma|=\int_a^b|\gamma'(t)|\,dt.
\]

Let $f\colon U\to\mathbb C$ be continuous. Let $F\colon U\to\mathbb C$ be holomorphic and satisfy $F'=f$. If $\gamma\colon[a,b]\to U$ is rectifiable, then
\[
\int_\gamma f(z)\,dz
=F(\gamma(b))-F(\gamma(a)).
\]


\paragraph{Proof (exercise; sketch).}
Establish
\[
\left|
\int_{\gamma|_{[a,t]}}f(z)\,dz
-\bigl(F(\gamma(t))-F(\gamma(a))\bigr)
\right|
\leq\varepsilon\left|\gamma\big|_{[a,t]}\right|
\]
for every $t\in[a,b]$ and every $\varepsilon>0$.

Notice that the common thread connecting
\[
F(\gamma(t))-F(\gamma(a))
\qquad\text{and}\qquad
\int_{\gamma|_{[a,t]}}f(z)\,dz
\]
is \(f(\gamma(a))(\gamma(t)-\gamma(a))\)
when $a$ and $t$ are close enough. This is the standard linear-approximation technique.
% BLOG-CONTENT-END
\end{document}
